\documentclass[11pt]{article}

\usepackage{amsmath,amssymb}
\usepackage{xurl}
\usepackage[colorlinks=true,linkcolor=blue,citecolor=blue,urlcolor=blue]{hyperref}

% Shared commands for notes in docs/paper-gaps/.

\DeclareUnicodeCharacter{2080}{\ifmmode{}_{0}\else\textsubscript{0}\fi}
\DeclareUnicodeCharacter{2081}{\ifmmode{}_{1}\else\textsubscript{1}\fi}
\DeclareUnicodeCharacter{2082}{\ifmmode{}_{2}\else\textsubscript{2}\fi}
\DeclareUnicodeCharacter{2099}{\ifmmode{}_{n}\else\textsubscript{n}\fi}

\newcommand{\Lean}{\textsc{Lean}}
\ExplSyntaxOn
\NewDocumentCommand{\leanid}{m}{
  \ifmmode
    \textsf{\detokenize{#1}}
  \else
    \group_begin:
    \urlstyle{sf}
    \tl_set:Nn \l_tmpa_tl {#1}
    \tl_replace_all:Nnn \l_tmpa_tl {\_} {_}
    \exp_args:NV \path \l_tmpa_tl
    \group_end:
  \fi
}
\ExplSyntaxOff
\providecommand{\tr}{\operatorname{tr}}
\newcommand{\tnleanIssueBase}{https://github.com/LionSR/TNLean/issues/}
\newcommand{\tnleanPrBase}{https://github.com/LionSR/TNLean/pull/}
\newcommand{\tnleanDocBase}{https://sirui-lu.com/TNLean/docs/}
\newcommand{\ghissue}[1]{\href{\tnleanIssueBase#1}{GitHub issue \##1}}
\newcommand{\ghpr}[1]{\href{\tnleanPrBase#1}{PR \##1}}
\newcommand{\doclink}[1]{\href{\tnleanDocBase#1.html}{\path{#1}}}

% Dirac notation and the identity operator, as in blueprint/src/macros/common.tex.
% \providecommand keeps note-local definitions authoritative.
\usepackage{dsfont}
\providecommand{\ket}[1]{|#1\rangle}
\providecommand{\bra}[1]{\langle #1|}
\providecommand{\braket}[2]{\langle #1 | #2 \rangle}
\providecommand{\ketbra}[2]{|#1\rangle\!\langle #2|}
\providecommand{\Id}{\mathds{1}}
\providecommand{\one}{\mathds{1}}
\providecommand{\C}{\mathbb{C}}
\providecommand{\MN}[1]{M_{#1}(\C)}
\providecommand{\MD}{M_D(\C)}

% Verdict marker: \gapnote{<kind>}{<status>}. Kind is the policy
% classification (clarification, local-correction, scope-restriction,
% unfaithful, false-source, open-gap); status is open, wip (elimination
% underway), resolved, or historical. Machine-read by the site index;
% prints a verdict line.
\newcommand{\gapnote}[2]{\par\noindent\textbf{Verdict:} #1 (#2).\par}


\title{PGWSVC08 Physical String Order versus Virtual-Boundary Nondecay}
\author{}
\date{2026-10-04}

\begin{document}
\maketitle
\gapnote{scope-restriction}{open}

\section{The source statement}

Let $\Psi$ be a translationally invariant state. P\'erez-Garc\'ia, Wolf, Sanz,
Verstraete, and Cirac define string order by the existence of a local unitary
$u\neq\mathbf{1}$ and local endpoint operators $x,y$, which may be taken
Hermitian, such that
\begin{align}
    S_N(x,y,u,\Psi)
        &:= \bra{\Psi}
            x\otimes u^{\otimes N}\otimes y
            \ket{\Psi},
    \label{eq:pgwsvc08_string_order_virtual_boundary_physical_correlator}\\
    \lim_{N\to\infty}
        \left|S_N(x,y,u,\Psi)\right|
        &>0.
    \label{eq:pgwsvc08_string_order_virtual_boundary_positive_limit}
\end{align}
The source denotes this definition by \path{SOP}
\cite[lines~112--122]{PerezGarcia2008StringOrder}.\footnote{Local source file:
\path{Papers/0802.0447/StringOrder-v10.tex}. A source footnote in
lines~124--127 cites Verstraete--Mart{\'i}n-Delgado--Cirac,
Phys.\ Rev.\ Lett.\ 92, 087201, for an ``alternative but equivalent''
definition in which the endpoints access the virtual levels. The source does
not prove this equivalence.}

The twist $u$ and the endpoint operators $x,y$ are all existentially
quantified. The fixed-endpoint order parameter used historically for the
Haldane chain is a special case.

For a pure finitely correlated state $\Psi_\infty$, let $\Lambda$ be its
positive virtual density matrix, let $\mathcal{E}_u$ be the transfer map
twisted by $u$, and let $\mathcal{E}_x,\mathcal{E}_y$ be the transfer maps
twisted by the endpoint operators. The source gives
\begin{align}
    S_N(x,y,u,\Psi_\infty)
        &=
        \tr\!\left[\Lambda
            \mathcal{E}_x
            \mathcal{E}_u^N
            \mathcal{E}_y(\mathbf{1})\right].
    \label{eq:pgwsvc08_string_order_virtual_boundary_source_transfer_form}
\end{align}
This is the source formula \path{SOPMP}
\cite[lines~176--181]{PerezGarcia2008StringOrder}.

Theorem~1 of Ref.~\cite{PerezGarcia2008StringOrder} gives an equivalent
criterion. A pure finitely correlated state has string order if and only if
there exist a unitary $\widetilde{u}\neq\mathbf{1}$, a virtual matrix $V$, and
physical indices $n,m$ such that
\begin{align}
    \mathcal{E}_{\widetilde{u}}(V)
        &=V,
    \label{eq:pgwsvc08_string_order_virtual_boundary_source_fixed_point}\\
    \tr\!\left(V\Lambda A_nA_m^\dagger\right)
        &\neq0.
    \label{eq:pgwsvc08_string_order_virtual_boundary_endpoint_realizability}
\end{align}
The source labels this result \path{Th1}
\cite[lines~270--276]{PerezGarcia2008StringOrder}. Condition
(\ref{eq:pgwsvc08_string_order_virtual_boundary_endpoint_realizability})
realizes the endpoint through $x=\ketbra{n}{m}$ and
$y=x^\dagger\widetilde{u}$, as shown in lines~257--268 of the same source.

The source removes
(\ref{eq:pgwsvc08_string_order_virtual_boundary_endpoint_realizability}) when
$x$ and $y$ may be products of observables on $D^2$ spins
\cite[lines~278--296]{PerezGarcia2008StringOrder}. For bond dimension $D$,
define
\begin{align}
    S_\ell
        &:=
        \operatorname{span}
        \left\{A_{n_1}\cdots A_{n_\ell}
            A_{m_\ell}^\dagger\cdots A_{m_1}^\dagger\right\}.
    \label{eq:pgwsvc08_string_order_virtual_boundary_endpoint_span}
\end{align}
For an injective tensor, these spaces grow strictly until they exhaust
$M_D(\mathbb{C})$. Endpoint operators on sufficiently many physical sites can
therefore realize every $D\times D$ virtual matrix.

\section{The restored peripheral phase}

In lines~241--255 of Ref.~\cite{PerezGarcia2008StringOrder}, choose right and
left peripheral matrices $V,Y$ for an eigenvalue
$\lambda=e^{i\theta}$. Iterating the twisted transfer map gives
\begin{align}
    \mathcal{E}_u^N(V)
        &=\lambda^N V,
    \label{eq:pgwsvc08_string_order_virtual_boundary_peripheral_iteration}\\
        &=e^{iN\theta}V.
    \label{eq:pgwsvc08_string_order_virtual_boundary_peripheral_phase}
\end{align}
The leading asymptotic contribution is consequently
\begin{align}
    S_N(x,y,u,\Psi_\infty)
        &=
        e^{iN\theta}
        \tr\!\left[Y\mathcal{E}_y(\mathbf{1})\right]
        \tr\!\left[\Lambda\mathcal{E}_x(V)\right]
        +o(1).
    \label{eq:pgwsvc08_string_order_virtual_boundary_corrected_asymptotic}
\end{align}
Lines~249--250 of the source print only the product of the traces and omit
$e^{iN\theta}$. This factor is required before the rephasing
$\widetilde{u}=e^{-i\theta}u$. The omission does not alter the absolute-value
criterion because
\begin{align}
    \left|e^{iN\theta}\right|
        &=1,
    \label{eq:pgwsvc08_string_order_virtual_boundary_phase_modulus}
\end{align}
but it changes the complex asymptotic formula.

The blueprint restores the phase in the proof of
\path{thm:pgwsvc08_physical_string_order} in
\path{blueprint/src/chapter/ch12_symmetry_string_order.tex}. The correlator is
formalized as \leanid{MPSTensor.physicalStringOrderParam}. The source theorem
and asymptotic argument remain marked \path{notready}; no \Lean{} theorem is
claimed for
(\ref{eq:pgwsvc08_string_order_virtual_boundary_corrected_asymptotic}). The
inline \path{Local fix} comment at the corrected display records this local
source correction.

\section{Formalized physical and virtual-boundary forms}

The transfer-form physical correlator is formally defined by
\begin{align}
    S_N(x,y,u)
        &:=
        \tr\!\left[\Lambda
            \mathcal{E}_x
            \mathcal{E}_u^N
            \mathcal{E}_y(\mathbf{1})\right].
    \label{eq:pgwsvc08_string_order_virtual_boundary_formal_physical_parameter}
\end{align}
\footnote{Formal declaration:
\leanid{MPSTensor.physicalStringOrderParam} in
\path{TNLean/MPS/Symmetry/StringOrderDefs.lean}.}
The source-style predicate is also formalized. It asserts that there exist a
unitary $u\neq\mathbf{1}$, physical endpoint operators $x,y$, and a number
$s>0$ such that
\begin{align}
    \left|S_N(x,y,u)\right|
        &\xrightarrow[N\to\infty]{}s.
    \label{eq:pgwsvc08_string_order_virtual_boundary_formal_positive_limit}
\end{align}
\footnote{Formal declaration:
\leanid{MPSTensor.HasPhysicalStringOrder} in
\path{TNLean/MPS/Symmetry/StringOrderDefs.lean}.}

The main formal theorems instead use virtual-boundary nondecay. For virtual
matrices $X,Y\in M_D(\mathbb{C})$ and a fixed twist $u$, define
\begin{align}
    T_L(X,Y,u)
        &:=
        \tr\!\left(\Lambda X\mathcal{E}_u^L(Y)\right).
    \label{eq:pgwsvc08_string_order_virtual_boundary_virtual_functional}
\end{align}
\footnote{Formal declaration:
\leanid{MPSTensor.stringOrderBoundaryParam} in
\path{TNLean/MPS/Symmetry/StringOrderDefs.lean}.}
Virtual-boundary nondecay means that there exist $X,Y$ and $c>0$ such that
\begin{align}
    \left\lVert T_L(X,Y,u)\right\rVert
        &\geq c
    \label{eq:pgwsvc08_string_order_virtual_boundary_uniform_nondecay}
\end{align}
for every length $L$.\footnote{Formal declaration:
\leanid{MPSTensor.HasStringOrder} in
\path{TNLean/MPS/Symmetry/StringOrderDefs.lean}.}

The formal development proves that virtual-boundary nondecay is equivalent to
virtual local covariance. It also proves virtual-boundary nondecay for
injective on-site-symmetric tensors satisfying the canonical normalization
hypotheses, together with the corresponding equivalence theorem.\footnote{
Formal declarations:
\leanid{MPSTensor.stringOrder_iff_localSymmetry},
\leanid{MPSTensor.hasStringOrder_of_symmetric_injective}, and
\leanid{MPSTensor.hasStringOrder_iff_of_symmetric_injective} in
\path{TNLean/MPS/Symmetry/StringOrder.lean}.}

The physical and virtual-boundary notions differ in four respects.

\begin{enumerate}
    \item
          \emph{Physical and virtual endpoints.}
          The source quantifies over physical endpoint operators $x,y$ with
          bounded site support. The formal virtual predicate quantifies over
          arbitrary matrices $X,Y$. A physical endpoint produces a particular
          virtual matrix in the transfer picture. Conversely, the span
          argument
          (\ref{eq:pgwsvc08_string_order_virtual_boundary_endpoint_span})
          realizes arbitrary virtual matrices for injective tensors, but this
          correspondence has not been formalized.

    \item
          \emph{Positive limit and uniform lower bound.}
          The source requires
          (\ref{eq:pgwsvc08_string_order_virtual_boundary_positive_limit}),
          whereas the virtual definition requires
          (\ref{eq:pgwsvc08_string_order_virtual_boundary_uniform_nondecay})
          for every length. A sequence with a positive limit may vanish at
          finitely many small lengths, so a fixed witness pair need not satisfy
          a uniform positive bound.

    \item
          \emph{Left-boundary form.}
          The physical transfer expression applies the endpoint through the
          adjoint channel, $\mathcal{E}_x^\dagger(\Lambda)$, whereas the formal
          virtual expression multiplies by $\Lambda X$. These forms have the
          same expressive range when $\Lambda$ is positive definite, as it is
          in the relevant theorems, but
          (\ref{eq:pgwsvc08_string_order_virtual_boundary_virtual_functional})
          is not the literal physical transfer expression.

    \item
          \emph{Fixed and existential twists.}
          Physical string order quantifies existentially over $u$. Virtual
          nondecay takes $u$ as a parameter, following the fixed-twist
          perspective of Theorem~2 in
          Ref.~\cite{PerezGarcia2008StringOrder}. The fixed-twist statement is
          finer, and existential quantification can be recovered from it.
\end{enumerate}

\section{Explicit physical-endpoint examples}

The two printed one-dimensional examples are now proved directly with
physical endpoint operators. These calculations use the stationary
infinite-chain transfer expression above; they do not identify it with an
arbitrary finite periodic-ring expectation. They also do not invoke the
one-site injectivity hypothesis of the general virtual-boundary criteria.

\paragraph{AKLT.}
In the physical order $(m=0,+1,-1)$, the printed tensor has letters
$(\sigma_z,\sqrt2\sigma_+,\sqrt2\sigma_-)/\sqrt3$. It is a diagonal unitary
physical rephasing of the existing AKLT tensor, which has a negative last
letter. The rephasing commutes with the chosen endpoints
$x=y=S_z=\operatorname{diag}(0,1,-1)$ and twist
$u=e^{i\pi S_z}=\operatorname{diag}(1,-1,-1)$.
\leanid{MPSTensor.akltPGWSVC08_sigmaZ_intertwine} exhibits the printed virtual
gauge $V=\sigma_z$, and \leanid{MPSTensor.akltSpinRotationZ_eq_exp} identifies
the actual matrix exponential.

The canonical density is $\Lambda=\Id/2$. It is positive, normalized and
stationary, as proved by \leanid{MPSTensor.akltPGWSVC08_canonical}.
The source's display \texttt{fixed} requires $\tr\Lambda=1$, whereas
Example~1 prints $\Lambda=\Id$. The latter is a normalization error:
\leanid{MPSTensor.physicalStringOrderParam_akltPGWSVC08} proves the printed
value $-4/9$ for every middle length $N\geq0$ with $\Id/2$, and
\leanid{MPSTensor.physicalStringOrderParam_akltPGWSVC08_identity_boundary}
proves that the literal identity boundary would give $-8/9$.
The concrete endpoints satisfy the physical positive-limit predicate,
\leanid{MPSTensor.akltPGWSVC08_hasPhysicalStringOrder}.

\paragraph{Cluster state.}
The source matrices are exactly the existing review tensor
\leanid{MPSTensor.clusterTensorRMP}, rather than its transposed representative.
With $u=-\sigma_x$, $V=\sigma_y$ and $\Lambda=\Id/2$, every endpoint trace
$\tr(V\Lambda A_nA_j^\dagger)$ vanishes. Moreover,
\leanid{MPSTensor.clusterTensorRMP_physicalStringOrderParam_eq_zero} proves
$S_N(x,y,-\sigma_x)=0$ for arbitrary one-site endpoints whenever $N\geq2$.
The absence assertion fixes this twist; it does not quantify over other
possible physical unitaries.

For the printed two-site endpoints $x=\sigma_z\otimes\sigma_y$ and
$y=\sigma_y\otimes\sigma_z$,
\leanid{MPSTensor.clusterTensorRMP_twoSite_physicalStringOrderParam} gives
\begin{align}
    \tr\!\left[\frac{\Id}{2}\,
        \mathcal E_x^{(2)}\mathcal E_{-\sigma_x}^{N}
    \mathcal E_y^{(2)}(\Id)\right] & =1
\end{align}
for every \emph{single-site} middle length, including odd lengths.
The blocked-index convention fixes the endpoint order. The corresponding
physical string-order predicate is also proved on the two-site blocked tensor
by \leanid{MPSTensor.clusterBlockedRMP_hasPhysicalStringOrder}, where the
middle length counts blocks and the twist is $u\otimes u$.

These results supply the printed data in Examples~1--2 of
\cite[lines~392--411]{PerezGarcia2008StringOrder}, tracked by \ghissue{8268}.
They do not close the general physical-endpoint realization or physical
reduced-density-operator symmetry theorems in \ghissue{4872}.

\section{Verdict}

Virtual-boundary nondecay is mathematically sound. For the injective tensors
covered by the relevant theorems, the existential physical-endpoint criterion
and the existential virtual-boundary criterion are equivalent by the source's
span argument
(\ref{eq:pgwsvc08_string_order_virtual_boundary_endpoint_span}). The physical
correlator and the source's positive-limit predicate are now formal
definitions, but the equivalence between physical and virtual endpoints has no
formal counterpart. Current general equivalence theorems therefore establish only the
virtual-boundary formulation; the explicit physical examples above are
proved independently.

There is a further hypothesis restriction in the general equivalence
results: their one-site injectivity assumption is stronger than the source's
purity condition (a positive stationary density and a unique peripheral
transfer eigenvalue). The explicit AKLT and cluster computations above do
not remove that restriction from Lemma~1 or Theorems~1--2.

The endpoint-realization restriction requires a theorem that realizes arbitrary
virtual boundary matrices by physical endpoint operators for injective
tensors, followed by a source-faithful proof of Theorem~1 in
Ref.~\cite{PerezGarcia2008StringOrder}. The existing single-site construction
\leanid{MPSTensor.physRealize} is a starting point.\footnote{Source path:
\path{TNLean/MPS/Chain/VirtualInsertion.lean}. This file is not currently
imported by the string-order files.} Separately, removing the purity-versus-injectivity restriction requires
a spectral and symmetry argument under the source canonical-purity
hypotheses; physical endpoint realization for an injective tensor does not
establish that generalization. Both tasks are tracked in \ghissue{4872};
\ghissue{2660} was absorbed into that issue.

\bibliographystyle{alpha}
\bibliography{references}

\end{document}
